\chapter{Marco teórico y estado del arte}
\label{cap:marco-teorico}

En este capítulo se presenta el estado del arte sobre el que se apoya el
resto del trabajo. Se aborda el catálogo de riesgos OWASP Top 10, las metodologías
habituales de seguridad ofensiva --- en particular la OWASP Web Security
Testing Guide, tomada como metodología de referencia a partir de la cual se
derivará en el capítulo~\ref{cap:metodologia} una metodología propia más
ágil --- y el estándar de puntuación de severidad
CVSS v3.1, el cual será empleado en el capítulo~\ref{cap:explotacion} para valorar cada
una de las vulnerabilidades explotadas.

\section{OWASP Top 10 y vulnerabilidades web comunes}

El OWASP Top 10~\citelib{OWASP21} es el catálogo de referencia del sector para
clasificar los riesgos de seguridad más críticos en aplicaciones web,
elaborado y actualizado periódicamente por la OWASP Foundation a partir de
datos proporcionados por organizaciones y profesionales de todo el mundo. Las seis
vulnerabilidades reproducidas en \breachlab corresponden a categorías de la
edición 2021 de este catálogo:

\begin{description}
\item[Inyección (A03).] Ocurre cuando datos no confiables se envían a
  un intérprete (SQL, sistema operativo, LDAP, etc.) como parte de un
  comando o consulta, con el objetivo de alterar su lógica original. El caso
  analizado dentro de \breachlab es una inyección SQL clásica en una consulta 
  de autenticación (CWE-89).

\item[Cross-Site Scripting — XSS (A03).] Permite a un atacante inyectar
  script en páginas vistas por otros usuarios, ejecutándose en el contexto
  del navegador de la víctima. En aplicaciones SPA modernas se manifiesta
  habitualmente como XSS basado en DOM: el dato malicioso nunca pasa por una
  plantilla del servidor, sino que se inserta directamente en el árbol DOM
  desde JavaScript (CWE-79).

\item[Cross-Site Request Forgery — CSRF (A01).] Dicha vulnerabilidad obliga 
  al navegador de una víctima autenticada a enviar, sin su conocimiento, una 
  petición a una aplicación en la que tiene sesión activa, aprovechando que el 
  navegador adjunta automáticamente las cookies de sesión (CWE-352).

\needspace{15\baselineskip}
\item[Fallos de identificación y autenticación (A07).] Se incluyen desde el
  almacenamiento de contraseñas en texto plano hasta
  la ausencia de límites de intentos, pasando por mensajes que permiten enumerar
  usuarios válidos o una gestión débil de la sesión (CWE-287, CWE-307).

\item[Control de acceso roto (A01).] Se produce cuando la aplicación no
  verifica de forma consistente que el usuario autenticado tenga permiso
  para operar sobre el recurso al que se desea acceder. El caso particular en el que un
  identificador de objeto se puede manipular para acceder a datos de otro
  usuario se conoce como IDOR, \emph{Insecure Direct Object Reference}
  (CWE-639).

\item[Configuración de seguridad incorrecta (A05).] Se incluye cualquier ausencia de
  \emph{hardening} (cabeceras de seguridad no configuradas, servicios con
  configuración por defecto, mensajes de error excesivamente detallados,
  etc.) que rara vez son explotables de forma aislada, pero que resultan en
  la reducción de la profundidad de defensa frente a otras posibles 
  vulnerabilidades (CWE-16).
\end{description}

Todos los identificadores CWE previamente citados corresponden al catálogo \emph{Common
Weakness Enumeration}~\citelib{MITRE}, mantenido por MITRE, que clasifica
tipos genéricos de debilidades de software con independencia de una
vulnerabilidad concreta. La tabla~\ref{tab:owasp-mapeo} resume la
correspondencia entre las seis vulnerabilidades a analizar dentro de \breachlab, con su categoría
en el OWASP Top 10 y su CWE asociado.

\begin{table}[htbp]
  \centering
  \begin{tabular}{lp{5.2cm}l}
    \toprule
    \tabhead{Vulnerabilidad} & \tabhead{Categoría OWASP} & \tabhead{CWE} \\
    \midrule
    Inyección SQL                        & A03 -- Inyección                              & CWE-89  \\
    XSS (DOM)                            & A03 -- Inyección                              & CWE-79  \\
    CSRF                                 & A01 -- Pérdida de control de acceso           & CWE-352 \\
    Fallos de autenticación              & A07 -- Fallos de identificación y autenticación & CWE-287, CWE-307 \\
    Control de acceso roto (IDOR)        & A01 -- Pérdida de control de acceso           & CWE-639 \\
    Configuración de seguridad insegura  & A05 -- Configuración de seguridad incorrecta  & CWE-16  \\
    \bottomrule
  \end{tabular}
  \caption{Relación de vulnerabilidades de \breachlab, 
    OWASP Top 10 y el CWE asociado.}
  \label{tab:owasp-mapeo}
\end{table}

\section{Metodologías de seguridad ofensiva}

La identificación y explotación de vulnerabilidades en este trabajo sigue el
enfoque de las pruebas de penetración de caja blanca: se dispone de acceso
completo al código fuente de la aplicación auditada, lo que permite localizar
con precisión el punto exacto en el que se introduce cada fallo y contrastar
el resultado de la explotación con la causa real en el código, algo que en
una auditoría de caja negra solo podría inferirse indirectamente.

Como marco de referencia para estructurar las pruebas se ha tomado la OWASP
Web Security Testing Guide (WSTG)~\citelink{OWASPWSTG}, que descompone la auditoría de una
aplicación web en fases: reconocimiento, gestión de la configuración, gestión
de identidad, autenticación, autorización, gestión de sesión, validación de
entradas, manejo de errores, etc. los cuales se recorrerán en el
capítulo~\ref{cap:metodologia} durante la definición de la metodología concreta aplicada a
\breachlab.

Como referencia complementaria sobre técnicas de identificación y
explotación, en particular para la inyección SQL y control de acceso, se
ha consultado también \emph{The Web Application Hacker's
Handbook}~\citelib{STU11}, obra de referencia en pentesting de aplicaciones
web.

Existen otros marcos de referencia habituales para estructurar una
auditoría de seguridad, como PTES~\citelink{PTES} (\emph{Penetration
Testing Execution Standard}), OSSTMM~\citelink{OSSTMM} (\emph{Open Source
Security Testing Methodology Manual}) o la guía técnica NIST~SP
800-115~\citelib{NIST80015}. A diferencia de la OWASP WSTG, estos
marcos se orientan a auditorías de infraestructura y de red de alcance más
general, y no ofrecen un catálogo de pruebas específico de aplicaciones web
tan detallado. Por esta razón se ha elegido la OWASP WSTG como metodología
de referencia de este trabajo: es la guía más extendida y detallada en el
ámbito concreto de la seguridad de aplicaciones web, y es también la
metodología de la que se parte en el capítulo~\ref{cap:metodologia} para
derivar una versión propia más ágil.

\section{CVSS v3.1 como marco de referencia}
\label{sec:cvss}

El \emph{Common Vulnerability Scoring System} (CVSS) es el estándar utilizado para expresar de forma cuantitativa y
comparable la severidad de una vulnerabilidad~\citelib{FIRST19}. La versión
3.1 calcula su
puntuación partiendo de ocho métricas: vector de ataque (AV), complejidad
de ataque (AC), privilegios requeridos (PR), interacción del usuario (UI),
alcance (S) y el impacto sobre confidencialidad (C), integridad (I) y
disponibilidad (A). El resultado es un valor entre 0.0 y 10.0 y una categoría
cualitativa (baja, media, alta o crítica).

En este trabajo, cada una de las seis vulnerabilidades del
capítulo~\ref{cap:explotacion} se acompaña de su vector CVSS v3.1 y de la puntuación obtenida al introducirlo
en la calculadora oficial de FIRST.org, razonando explícitamente el valor
asignado a cada métrica en función del escenario concreto de \breachlab.

En noviembre de 2023, FIRST.org publicó CVSS
v4.0~\citelib{FIRST23}, que incorpora nuevas métricas (requisitos de seguridad y de impacto en sistemas subsecuentes) y redefine el
tratamiento del alcance. Para el desarrollo de este trabajo se mantiene la versión 3.1
por ser, en el momento de redacción del mismo, la más extendida en herramientas y
bases de datos de vulnerabilidades de uso habitual, además de
resultar suficiente para expresar el impacto de las seis vulnerabilidades
analizadas en \breachlab.
